% Verbatim copy of /Users/alingalatan/workspaces/grothe/handoff/manuscript/legacy/ch19-spectral-review.tex (file dated 2026-10-25). Only this line was added.
% Consolidated source module (formerly Chapter 19). Owner: lower-bound/07-remote-spectral.tex
\section{Spectral obstruction: audit, quantitative scope, and extra charges (GPT review package \code{spectral-handoff-review}, 15 October 18:27)}
\label{ch:spectral-review}
\chapsummary{Artificial radial laws expose limits of the scalar moment/Stein relaxation. These are not counterexamples to cubic-chaos K. The spectral identities and quantitative obstruction are retained.}

\paragraph{Provenance (IDE note, 15 October, evening).} This chapter is lower-bound GPT's \code{tex\_proposed/ch19-spectral-review.tex} from
\pathref{\~{}/Documents/Codex/2026-10-10/plea/outputs/spectral-handoff-review-20261015/} (zip \code{conjecture-k-spectral-reviewed.zip}), ported verbatim except for this header and the closing
reading note; the review, both versions of the reviewed note, the standalone \code{spectral\_review.tex}, the patches and logs are archived in
\pathref{handoff/research-third-decimal/gpt-reviews/spectral-20261015/}. \emph{What is reviewed}: not one of our chapters but an external ``spectral Conjecture K handoff'' (built on our
bundle of 16:45) which (i) constructs artificial radial laws $p_b=e^{-u}(1-bL_4+b^2h)$ that satisfy all the scalar inputs of the lower certificate (the moment constraints of
\autoref{ch:certified}, $\E\sqrt U<\rho$, the fluctuation budget and \emph{every} radial Stein cut) yet beat the Gaussian payoff at $\sigma^*$, so exact K-a is not provable from that scalar
relaxation; (ii) excludes them by the Azmoodeh--Malicet--Mijoule--Poly spectral inequalities for third chaos; (iii) proposes a joint spectral hierarchy retaining the mixed content.
GPT's review repairs a false uniform density bound in (i), proves the optimised excess of the whole family is $<4.65\cdot10^{-6}$, confirms (ii)--(iii) and adds the exact remainder (S6)
and the charge-three block (S7). Our replay: all four shipped verifiers PASS under \code{python3 -I} from an isolated directory (the received checker also passes, since its endpoint test
misses the false uniform claim); own exact checks of $L_4(64)=536321$, the $y=u-32$ identity, the seven Bernstein coefficients, the rational bound $5325076801/1146880000000000$ and
$723420b+424/3\le692$; Monte Carlo of $M_6\ge3M_4^2/(2v)+2\|\Pi_{1,0}R\|^2$ on random $\Hc21(\C^2)$ tensors (slack far above noise). Log \pathref{kattack-2026-10-15/gpt\_spectral\_replay.log},
script \code{spectral\_own\_check.py}. Status labels: \st{P} GPT for the analytic statements, \st{C} for the exact finite checks (replayed), \st{F} for the quadrature remarks.
\textbf{No endpoint changes} (certified bound at the time $\KG>1.40160592$; the current Lean bound is $\KG>1.401674103\stale{22 Oct}{now $\KG>1.401695322458$}$, \S\ref{sec:certcombinede}).

\paragraph{Status.}
The received spectral handoff gives an obstruction to exact Gaussian
domination by a scalar relaxation, not a counterexample to Conjecture K.
Its uniform density estimate needs the repair below. The spectral identities
are valid. This chapter also proves an upper bound of $4.65\cdot10^{-6}$
on the optimized excess of the entire constructed artificial family, and
adds a charge-three matrix constraint using moments already present in the
proposed degree-two hierarchy. None of these results changes a numerical
Grothendieck endpoint or proves the required $K_\eta$ for general cubic chaos.
The analytic arguments are not newly formalized in Lean.

\subsection{Repair of the uniform density estimate}

Let $L_j$ be the ordinary Laguerre polynomial and set
\[
 h=L_4^2-1+8L_1-42L_2+152L_3,\qquad
 P_b=1-bL_4+b^2h,\qquad 0<b\le10^{-4}.
\]
Write $h=L_4^2+C_3$, where
\[
 C_3(u)=117-380u+207u^2-\frac{76}{3}u^3,
 \qquad C=\frac{3163103}{3}.
\]
The original proof claimed $P_b(u)\ge3/4-Cb^2$ globally. This is false:
at $u=64$ and $b=1/1072642=1/(2L_4(64))$,
\[
 L_4(64)=536321,\qquad C_3(64)=-\frac{17451937}{3}<-C,
\]
so $P_b(64)=3/4+b^2C_3(64)<3/4-Cb^2$.
The original verifier compared its two regional lower bounds only at
$b=10^{-4}$, which does not establish that comparison throughout the interval.

The replacement
\[
 \boxed{P_b(u)\ge D_b:=\frac34-106b\ge0.7394}
 \tag{S1}
\]
holds uniformly. On $[0,32]$, completing the square gives
$P_b\ge3/4-Cb^2\ge D_b$. For $u\ge32$, put $y=u-32$. The exact identity
\[
 L_4(u)-\frac{u^4}{48}
 =2945+\frac{2612}{3}y+67y^2+2y^3+\frac{y^4}{48}
\]
proves $L_4(u)\ge u^4/48$. Consequently
\[
 |C_3(u)|\le33u^3\le\frac{99}{2}L_4(u),
\]
and a second square completion yields
\[
 P_b(u)\ge\frac34-\frac{99}{4}b-\frac{9801}{16}b^2\ge D_b.
\]
For the density $p_b(u)=e^{-u}P_b(u)$, the original Laguerre calculation of
the scalar Stein kernel $w_b$ now gives the valid estimate
\[
 \operatorname{Var}(w_b)
 \le b^2\frac{1-692b+723420b^2}{3/4-106b}
 \le\frac43b^2.
 \tag{S2}
\]
The final inequality is precisely $723420b+424/3\le692$ and holds on the
whole stated interval. Thus the obstruction and its parameter range survive.
The corrected checker also tests the exact failure of the old denominator.

In the visible-profile argument, the assumptions should explicitly be
$\beta\ge0$ and $\ell-2\beta>0$. They imply that
$L=Z(\ell+\beta(|Z|^2-2))$ is a proper Gaussian variance mixture,
$L\mid V\sim\sqrt V W$. The peak parameters satisfy these assumptions.

\subsection{The obstruction has uniformly small optimized excess}

Define, for $\sigma>0$,
\[
 f_\sigma(u)=2\mathbb E|L+\sigma\sqrt u|,
 \qquad I_n(\sigma)=\int_0^\infty e^{-u}f_\sigma(u)L_n(u)\,du.
\]
For every integer $n\ge1$, the exact identity is
\[
 \boxed{-I_n(\sigma)=
 \frac{\Gamma(n-\tfrac12)}{2\Gamma(n+1)}\,
 \sigma\,\mathbb E\left[
 \left(\frac{\sigma^2}{V+\sigma^2}\right)^{n-1/2}\right].}
 \tag{S3}
\]
To prove it, use the Laplace representation of the square root and the
Gaussian integral conditionally on $V$. The Laguerre transform leaves
\[
 \frac{\sigma^{2n}}{\sqrt\pi}
 \int_0^\infty\frac{t^{n-3/2}}{(1+t(V+\sigma^2))^{n+1}}\,dt.
\]
The beta integral evaluates this expression as (S3). Subtracting the
constant term before the Laguerre integral, or using the integrability
condition of the Bernstein representation, justifies interchange.

Put $x=\sigma^2/(V+\sigma^2)$. Since
\[
 h=2L_2+346L_4-520L_5+500L_6-280L_7+70L_8,
\]
formula (S3) gives
\[
 \int_0^\infty e^{-u}f_\sigma(u)h(u)\,du
 =-\sqrt\pi\,\sigma\,\mathbb E[x^{3/2}P(x)],
\]
where
\[
 P(x)=\frac14+\frac{865}{64}x^2-\frac{455}{32}x^3
 +\frac{2625}{256}x^4-\frac{1155}{256}x^5
 +\frac{15015}{16384}x^6.
\]
The degree-six Bernstein coefficients of $P$ on $[0,1]$ are
\[
 \left(\frac14,\frac14,\frac{221}{192},\frac{287}{128},
 \frac{895}{256},\frac{7399}{1536},\frac{101671}{16384}\right).
\]
In particular, $P\ge1/4$, so the second-order payoff perturbation is
negative. The first-order coefficient satisfies
$-I_4\le(5\sqrt\pi/128)\sigma$. With $c=5\sqrt\pi/128$, therefore,
\[
 \mathbb E_b f_\sigma(U)
 \le\mathbb E_{\mathrm{Exp}(1)}f_\sigma(U)+cb\sigma.
 \tag{S4}
\]

Write
\[
 B(\sigma)=\sqrt\pi\,\mathbb E\sqrt{V+\sigma^2}
             -\frac{\sigma^2}{\varepsilon},\qquad
 S=\sup_{\sigma\ge0}B(\sigma).
\]
For $\sigma>0$,
$B'(\sigma)\le\sqrt\pi-2\sigma/\varepsilon$. Consequently the upper
envelope $B(\sigma)+cb\sigma$ is decreasing past
$R=\varepsilon(\sqrt\pi+cb)/2$, and (S4) proves
\[
 \boxed{\sup_{\sigma\ge0}
 \left[\mathbb E_b f_\sigma(U)-\frac{\sigma^2}{\varepsilon}\right]-S
 \le\varepsilon\frac{5\pi}{256}b\left(1+\frac{5b}{128}\right).}
 \tag{S5}
\]
At $\varepsilon=1891/2500$ and $b\le10^{-4}$, using $\pi<22/7$, this is at most
\[
 \frac{5325076801}{1146880000000000}<4.65\cdot10^{-6}.
\]
Thus this entire family is compatible with the useful $10^{-4}$ slack,
although it disproves exact domination in the old unrestricted relaxation.
Its kurtosis excess is at most $4b^2\le4\cdot10^{-8}$, outside the stated
application window $0.06\le\mathbb E U^2-2\le4$.
No bound on all other laws in the relaxation is asserted here.

\subsection{The retained content and its exact remainder}

Let $q\ne0$ belong to $\mathcal H_{2,1}$, let $Nq=3q$ for the positive
Gaussian number operator, and set
\[
 U=|q|^2,\quad v=\mathbb E U,\quad M_4=\mathbb E U^2,
 \quad M_6=\mathbb E U^3,\quad R=qU,\quad\mu=M_4/v.
\]
The received hierarchy implies
\[
 M_6\ge\frac{3M_4^2}{2v}+2\|\Pi_{1,0}R\|_2^2.
\]
The following exact identity identifies all discarded terms:
\begin{align*}
 M_6-\frac{3M_4^2}{2v}-2\|\Pi_{1,0}R\|_2^2
 ={}&\frac32\|\Pi_{2,1}R-\mu q\|_2^2
      +\|\Pi_{3,2}R\|_2^2+\frac12\|\Pi_{4,3}R\|_2^2\\
 &+\frac16\|\Pi_{3,0}(q^3)\|_2^2
      +\frac19\|\Pi_{4,1}(q^3)\|_2^2
      +\frac1{18}\|\Pi_{5,2}(q^3)\|_2^2.
 \tag{S6}
\end{align*}
Indeed, if $A=\mathbb E[U^2\Gamma(q,\bar q)]$, integration by parts gives
\begin{align*}
 \langle R-\mu q,(9-N)(R-\mu q)\rangle
    &=3M_6+A-6M_4^2/v,\\
 \langle q^3,(9-N)q^3\rangle&=9(M_6-A).
\end{align*}
Add one ninth of the second identity to the first, subtract
$8\|\Pi_{1,0}R\|_2^2$, and divide by four. The allowed spectral degrees
give (S6). This is dimension independent. It is an explanation of the
received scalar inequality, not a new constraint beyond its degree-one
hierarchy and upper bound on $a_2$.

\subsection{A stronger matrix using the same moment variables}

Now use the normalized notation
\[
 \mathbb E|q|^2=1,\quad w=\Gamma(q,\bar q)/3,
 \quad m_j=\mathbb E U^j,\quad a_j=\mathbb E[wU^j].
\]
The received charge-one hierarchy uses $qU^i$. At $i\le2$, its variables
already include $m_1,\ldots,m_5$ and $a_0,\ldots,a_4$. With no additional
moments one can impose
\[
 \boxed{\begin{pmatrix}
 15m_3-27a_2&12m_4-24a_3\\
 12m_4-24a_3&9m_5-21a_4
 \end{pmatrix}\succeq0.}
 \tag{S7}
\]
This is the quadratic form $15-N$ on $\{q^3,q^3U\}$, both of degree at most
15. The three energy entries are $27a_2$, $3m_4+24a_3$, and
$6m_5+21a_4$, respectively. This proves (S7).

For precision, (S7) is not implied by the displayed charge-one constraints
alone. An exact finite-matrix example is
\begin{align*}
 (m_0,\ldots,m_5)&=(1,1,4,1268/49,91296/343,8039120/2401),\\
 (a_0,\ldots,a_4)&=(1,2,300/49,7608/343,1607824/2401).
\end{align*}
Take $Q_1=xx^*$, $Q_4=yy^*$, all other $Q_r=0$, where
\[
 x=(1,4,1268/49),\qquad y=(0,22/7,2536/49).
\]
These satisfy the received sums, spectral supports, and $a_j$ bounds,
basic radial and $w$-weighted moment positivity, and the degree-four real
spectral matrix. But the determinant in (S7) is
$-412117447680/117649<0$.
This example is not asserted to be a realizable joint law or a cubic chaos.
The received report already recommends including the full real spectral
hierarchy; the example explains why it should not be dropped.

The general extension is explicit. For $k\ge1$, put $F_i=q^kU^i$,
$t=k+i+j-1$, and, when $t>0$,
\[
 h=\frac{k(i+j)+2ij}{t}.
\]
Its Gram and energy matrices are
\[
 G_{ij}=m_{k+i+j},\qquad
 D_{ij}=3\bigl[h\,m_{t+1}+(k^2-h)a_t\bigr].
 \tag{S8}
\]
The sole exception is $k=1,i=j=0$, with $D_{00}=3m_1$.
To verify (S8), apply the chain rule to $F_i,\overline{F_j}$ and use
\[
 m_{t+1}=(t+1)a_t+
 t\,\mathbb E[(\Gamma(q,q)/3)\bar q^{\,2}U^{t-1}].
\]
The last expectation is real by self-adjointness. The charge-$k$ spectral
blocks have eigenvalues $k+2r$, row $i$ supported on $r\le k+3i$;
their sums, unweighted and weighted by these eigenvalues, equal $G$ and $D$.

\paragraph{Implementation limits.}
The degree-two charge-three block is a cheap strengthening of the proposed
relaxation. It does not prove the optimum has the desired value.
The actual profiles still need rigorous approximation and tail errors.
Finite polynomial spectral cutoffs cannot be applied directly to
nonpolynomial profiles. The older exact-profile annihilators use
\emph{unnormalized} $\Gamma$; conversion to the present $w$ requires the
factor three. None of these finite checks replaces the analytic proof or
the missing finite-amplitude mixed-sector estimates.

\paragraph{Source and replay.}
The spectral moment method is due to Azmoodeh, Malicet, Mijoule and Poly,
\emph{Generalization of the Nualart--Peccati criterion},
Annals of Probability 44 (2016), 924--954, Theorem 3.1,
\url{https://arxiv.org/pdf/1305.6579}.
The accompanying files \texttt{corrected/verify\_obstruction.py},
\texttt{check\_scope\_bound.py}, and
\texttt{check\_spectral\_extensions.py} run with isolated standard-library
Python. They check the repaired constants, positive Bernstein coefficients,
exact matrix example, and the Wick instance of (S6), whose remainder is 8192.

\subsection{IDE reading note}\label{sec:spectral-reading}
\emph{Relation to our chapters.} (a) \autoref{ch:certs}, Task~B, found that the Lean-proved moment constraints \emph{alone} give a negative deficit near $\sigma^*$ ($d=-6\cdot10^{-5}$ at
$x\to0$ in the HiGHS solve; the exact optimum is $d\ge-9.8\cdot10^{-6}$, \autoref{ch:rescub}). The reviewed note shows more: even after adding the fluctuation budget and the entire radial Stein-cut family, the scalar relaxation still admits laws above $S$, so exact K-a
($\eta=0$) cannot be certified from radial moments and cuts, however many; GPT's bound (S5) shows this particular family exceeds $S$ by at most $4.65\cdot10^{-6}$, well inside the $K_\eta$
budget, consistent with the ``genuine moment-only loss $\approx10^{-5}$'' recorded in the K report. (b) Nothing here contradicts \autoref{ch:kjoint}--\autoref{ch:shared-residual-review}:
the certificates use the cuts as valid inequalities with their own slack and never claim exact domination; the exact-profile annihilators of \autoref{ch:crossbound} use the unnormalised
$\Gamma$ while this chapter uses $w=\Gamma(q,\bar q)/3$ (a factor $3$ to carry when combining them). (c) \emph{Lean relevance.} No Lean-proved statement is touched: the lower theorem assumes
\code{Moments} and \code{SteinInputs} as hypotheses for the abstract law (since commit \code{2c54991} inside \code{SharedReduction}\stale{22 Oct}{proved in Lean: \code{sharedReduction\_frame}, \code{reductionHyp\_of\_nodeCaps}; only \code{NodeCaps} remains}) and proves $V\le B$; the artificial laws satisfy those hypotheses and the conclusion is still true for them (the
certificate's slack absorbs their $<4.65\cdot10^{-6}$), so there is no error anywhere. No hypothesis changes. \emph{Possibly sharper inputs:} the spectral inequalities (9), (10) of the reviewed
note and (S7) here are rigorous necessary conditions for third chaos (AMMP, Theorem~3.1, \st{P} literature) and could be added to the residual layer (\code{Moments}) as further cuts, which
would exclude exactly the laws that make Task~B's deficit negative; the received note's caveat applies (lower moment bounds alone need not improve an unbounded-support optimisation without
upper tail information), and no quantified gain exists yet, so this is a candidate, not a result. (d) The joint spectral hierarchy of the reviewed note is a numerical research proposal for
(T$_0$) (\autoref{ch:weighted}); it needs certified profile approximation and tail errors, and is complementary to the exact-profile routes of \autoref{ch:kjoint} and \autoref{ch:crossbound}.
